% defined-in: zorich (page 131)
% === SEMANTIC MAPPING ===
% 1. Variables & Types:
%     - f: Function (f \colon X \rightarrow \RealNumbers)
%     - X: Set (The domain of the function, typically a subset of \RealNumbers)
%     - x_0: Point (The point where the limit is evaluated, x_0 \colon X)
%     - L: Real Number (The value of the limit, L \colon \RealNumbers)
%     - V: Neighborhood (A set in the codomain \RealNumbers)
%     - B_prime: Base element (A set in the domain X, representing the approach from the left)
% 2. Data Structures:
%     - The concept of 'left approach' is handled by the specialized macro \mathrm{IsOneSidedLimit}(f, x_0, L), which implicitly restricts the base elements B_prime to those approaching x_0 from below.
%     - Sets (Domains and Codomains) are represented using the \Function and \RealNumbers macros.
% 3. Implicit Bounds:
%     - The definition requires the existence of a domain X (often a subset of \RealNumbers) and a point x_0 \colon X.
%     - The definition relies on the standard topological definition of neighborhoods \Neighborhood.
% ========================

% entity-id: def-real-analysis-limit-one-sided
% entity-type: def
% macro: DefinitionOneSidedLimit
% args: 3
% notation: \mathrm{IsOneSidedLimit}(f, x_0, L)
\hypertarget{def-real-analysis-limit-one-sided}{}
\begin{definition}[Left Limit of a Function]
\TerForall X \colon \Set \TerForall f \colon \Function{X}{\RealNumbers} \TerForall x_0 \colon X \TerForall L \colon \RealNumbers \\
\mathrm{IsOneSidedLimit}(f, x_0, L) \coloneqq \left( L = \lim_{x \to x_0^-} f(x) \right)
\end{definition}